\section{Введение}
\label{sec:Chapter0} \index{Chapter0}

    В настоящее время создается множесто язуков программирования, каждый из которых своей целью ставит 
удобство решения определенного класса задач. Большинство ныне существующих языков программирования
относятстя к платформонезависимыми языкам. Главным их преимуществом является высокая переносимость
кода, существует идиома описывающая данный вид поведения(WORA — Write Once, Run Anywhere). Примерами
таких языков являются: ArkTS, C\# Dart, Java. 

Высокая переносимость обеспечивается по средствам специальных сред исполнения, называемых 
виртуальными машинами, которые исполняют байт-код. Байт-код -- промежуточный, платформенно-независимый 
код низкого уровня, получаемый при компиляции исходного текста программы. Он состоит из компактных 
инструкций, которые исполняются не напрямую процессором, а виртуальной машиной 
(например, JVM для Java), что позволяет запускать приложения на Windows, Linux, macOS и других 
платформах без перекомпиляции.

В связи со сказанным выше на первый план выходит проблемя отслеживания корректности(здесь и далее под
корректностью будем понимать соответсвие архитектуре набора команд и семантике языка), 
сгенерированного байт-кода. Программы отслеживающий корректность назывются верификаторами байт-кода.
Важность изучения и разработки данных программы обуславливается тем, что исполняемые приложения 
могут отвечать за функционирование
крайне чувствительных областей, таких как телекоммуникации, банковская сфера, медицинские услуги, 
где ошибка исполнения может привести к необратимому ущербу. Наличие 
верификатора позволяет обнаруживать ошибки до времени исполнения. Кроме того, верификатор 
оказывается полезен на этапе разработки языка программирования.

Существует несколько источников некорректного байт-кода:
\begin{enumerate}
    \item Ошибка компилятора на этапе кодогенерации, Lowering(понижение), регистровой аллокации и т.д.
    \item Повреждения пакета с исполняемым кодом во время его передачи по сети.
\end{enumerate}

------

В данной работе фокус идет на повышение точности работы верификатора байт-кода путем расширения его 
типовой система и увеличения множества верифицируемых инструкций.